% =========================================================
% setup/00_metadata.tex
% Public, user-editable portfolio metadata
% =========================================================
%
% This is the first file to edit when starting a portfolio project.
% Keep plain-text PDF fields free of manual line breaks and commands.
% Escape LaTeX special characters when needed: & -> \&, % -> \%, _ -> \_.

% -----------------------------------------------------------------
% 1. Public project identity
% -----------------------------------------------------------------
\newcommand{\projecttitle}{Compact Sensor Array Design and Verification}
\newcommand{\projectsubtitle}{Analytical Development, Numerical Modeling, and Engineering Evidence}
\newcommand{\projectdomain}{Applied Electromagnetics and Embedded Sensing}
\newcommand{\projectshorttitle}{Compact Sensor Array}

% Formatted title for the cover.
\newcommand{\reporttitle}{%
    Compact Sensor Array Design\\[0.18cm]
    and Verification%
}

% Plain title for PDF metadata and accessibility.
\newcommand{\reporttitleplain}{Compact Sensor Array Design and Verification}

% -----------------------------------------------------------------
% 2. Public author identity
% -----------------------------------------------------------------
\newcommand{\authorname}{Author Name}
\newcommand{\authorrole}{Engineer and Technical Researcher}
\newcommand{\portfolioedition}{Independent Engineering Portfolio Project}
\newcommand{\publicationdate}{\today}

% Public links. Replace placeholders before publication.
\newcommand{\repositoryurl}{https://github.com/username/project-repository}
\newcommand{\linkedinurl}{https://www.linkedin.com/in/username}
\newcommand{\contactemail}{name@example.com}

% Set each switch to true or false.
\newif\ifshowauthorrole
\showauthorroletrue
\newif\ifshowrepository
\showrepositorytrue
\newif\ifshowlinkedin
\showlinkedinfalse
\newif\ifshowcontactemail
\showcontactemailfalse

% -----------------------------------------------------------------
% 3. Optional public logo or project mark
% -----------------------------------------------------------------
\newif\ifshowprojectlogo
\showprojectlogofalse
\newcommand{\projectlogo}{figures/project-logo-placeholder.png}
\newcommand{\projectlogowidth}{4.6cm}

% -----------------------------------------------------------------
% 4. Portfolio context and evidence declarations
% -----------------------------------------------------------------
\newcommand{\projectorigin}{%
This public edition was prepared from an independently completed engineering project.
Replace this sentence with a concise, factual description of the project origin.}

\newcommand{\publiceditionnote}{%
Administrative, grading, student-identification, and institutional material has
been removed. The public edition focuses on technical decisions, evidence, limitations,
and reproducibility.}

\newcommand{\analyticalevidence}{%
Closed-form calculations, governing equations, design rules, and independent scripts.}

\newcommand{\simulatedevidence}{%
Numerical model construction, solver settings, parameter sweeps, and exported results.}

\newcommand{\measuredevidence}{%
No fabricated prototype or measurement campaign is claimed in this example template.}

\newcommand{\numericalscope}{%
Replace this paragraph with solver, mesh, boundary, material, convergence, and software
limitations that materially affect interpretation of the reported results.}

\newcommand{\affiliationdisclaimer}{%
This document is an independent portfolio artifact and does not imply endorsement by an
employer, university, instructor, customer, software vendor, or standards organization.}

% -----------------------------------------------------------------
% 5. Running headers and PDF properties
% -----------------------------------------------------------------
\newcommand{\headerleft}{Engineering Portfolio}
\newcommand{\headerright}{\projectshorttitle}
\newcommand{\pdfsubjecttext}{\projectdomain}
\newcommand{\pdfkeywordstext}{engineering portfolio, analytical design, numerical simulation, verification, reproducibility, technical documentation}
